\documentclass[11pt]{article}
\usepackage[margin=25mm]{geometry}
\usepackage{fontspec}
\setmainfont{Latin Modern Roman}
\usepackage{amsmath,amssymb,mathtools}
\usepackage{unicode-math}
\setmathfont{Latin Modern Math}
\usepackage{xurl,hyperref}
\hypersetup{colorlinks=true,urlcolor=blue}
\setcounter{secnumdepth}{0}
\setlength{\parindent}{0pt}
\setlength{\parskip}{5pt}
\setlength{\emergencystretch}{3em}
\newcommand{\sref}[2]{\hyperlink{source:#1:#2}{[S#2]}}
\newcommand{\cataloguescope}[1]{\paragraph{Recorded target.}#1}
\newcommand{\R}{\mathbb R}
\newcommand{\C}{\mathbb C}
\newcommand{\N}{\mathbb N}
\newcommand{\Z}{\mathbb Z}
\newcommand{\Q}{\mathbb Q}
\newcommand{\F}{\mathbb F}
\newcommand{\eps}{\varepsilon}
\begin{document}
\section{E0-Semigroups from Free Product Systems}\label{statement}

\subsection{Definitions and mathematical statement}

\paragraph{Shared notation.}$\mathbb N=\{1,2,\ldots\}$; the standard number fields have their usual meanings. All additional conventions are stated locally.


\paragraph{Definitions and conventions.} For a unital coefficient algebra $B$, a Hilbert $B$-bimodule is a right Hilbert $B$-module with a compatible left action. $\mathcal B^a(E)$ denotes its adjointable operators. A pointed free product system $(E_t,\omega_t)$ has central reference vectors and associative identifications $E_s*E_t\simeq E_{s+t}$, using the source's reduced free product over $B$. A tensor product subsystem instead has $E_s\odot E_t\simeq E_{s+t}$ with $\odot$ the internal tensor product over $B$. An $E_0$-semigroup is a continuous semigroup of unital $*$-endomorphisms. A unit is a section whose values multiply according to the product-system structure. The reference vectors define an inductive limit $E_\infty$ and a free left dilation $v_t:E_\infty*E_t\to E_\infty$. The source uses this dilation and the induced unital embeddings to construct an $E_0$-semigroup on $\mathcal B^a(E_\infty)$. A noise here has the reference-vector/vacuum structure specified in that construction.

\paragraph{Source scope and qualifications.} The free left dilation is already constructed in the report. The unresolved target is which noises arise by this construction, not existence of the construction alone.

\paragraph{Catalogue mathematical target.}
\paragraph{Catalogue wording retained for traceability.}
Which $E_0$-semigroups (necessarily noises) arise from the free-product-system construction and its associated free left dilation?

\paragraph{Saved literature context.} The supplied catalogue's assessment dated 2026-08-21 records: The broad characterization of E0-semigroups arising from free product systems remains an open classification problem; existing work constructs important families but not all cases. The assessment is retained as dated provenance; it is not a new comprehensive literature or open-status certification.
\subsection{Short English statement}
Characterize all endomorphism semigroups obtainable from the report's free product-system dilation.
\subsection{Research attempt}
\paragraph{Model, date and investigated scope.}
OpenAI GPT-6 Astra (configured model \texttt{gpt-6-astra}), ultra reasoning; generated 11 October 2026. Derives a fourth-moment obstruction to realizing the same nondegenerate increments as both tensor independent and free.

\paragraph{Sources and literature check.}
The complete supplied definition, including source qualifications and the dated literature assessment, was reviewed. The calculation below is an elementary restricted attempt motivated by [S2]; it does not claim novelty or a fresh literature-status determination.

\paragraph{Context and investigated scope.} The complete supplied definition, including source qualifications and the dated literature assessment, was reviewed. The calculation below is an elementary restricted attempt motivated by [S2]; it does not claim novelty or a fresh literature-status determination.

\paragraph{An increment-moment obstruction.} In the scalar-state case, suppose candidate past and future increment algebras $A_1,A_2$ are freely independent in a state $\phi$. By the defining centered alternating-word condition, self-adjoint $x\in A_1,y\in A_2$ with $\phi(x)=\phi(y)=0$ satisfy
\[
 \phi(xyxy)=0.
\]
By contrast, if these same increment algebras commute and their state is a tensor product, then $xyxy=x^2y^2$ and $\phi(xyxy)=\phi(x^2)\phi(y^2)$. With both variances positive the two requirements are incompatible. A concrete tensor example is $x=Z\otimes I,y=I\otimes Z$ in the normalized trace on $M_2\otimes M_2$: the alternating moment is one.

\paragraph{A degenerate overlap.} If one centered increment algebra is zero in its state representation, the obstruction disappears. This includes the one-dimensional pointed free system, whose dilation is the identity semigroup. Thus the moment test has the correct trivial endpoint and cannot be used to exclude all noises with a tensor presentation.

\paragraph{What classification still needs.} The calculation only tests whether a specified pair of increments can simultaneously be tensor independent and free under the same state. An endomorphism semigroup might admit another increment realization, another reference vector, or an equivalence which changes these choices. Therefore it is not an intrinsic classification invariant of the bare semigroup. The first gap is a representation-independent reconstruction criterion for the source's free dilation. The already constructed free left dilation is not reproved or presented as the unresolved result.

\paragraph{Verification and adversarial audit.}
Checked centered alternating moments, explicit Pauli example, zero-variance endpoint and dependence on the chosen increment/state realization. Model self-audit only; no independent human proof checking. Prompts A and B were applied to this canonical record and its qualified source scope. Checked centered alternating moments, explicit Pauli example, zero-variance endpoint and dependence on the chosen increment/state realization. Retained unresolved status and zero original-target progress; No intrinsic criterion reconstructing a free-product-system dilation from an arbitrary noise is established.

Coordinator checked the centered free alternating moment and incompatible positive tensor-product moment, with degenerate and representation-dependent cases qualified. This additional check was model-assisted, not independent human verification.

\paragraph{Tools, process and contributions.}
Python 3 was used for file preparation and structural checks; no numerical experiment or theorem prover was used. Contribution prompt A was followed by prompt B. GPT-6 Astra with ultra reasoning drafted the research and performed its stated self-audit; any additional coordinator contribution is identified in the verification and audit. A GPT-6-based Codex coordinator assembled the manuscript and checked repository integration; its exact variant and reasoning configuration were not exposed. No human mathematical contribution, independent human proof check or formal proof certification is supplied. Complete derivations and the consulted source references are included here; private reasoning transcripts are not published.

\paragraph{Final status and first remaining gap.}
\textbf{UNRESOLVED}. No intrinsic criterion reconstructing a free-product-system dilation from an arbitrary noise is established.

\section{Completion Estimate}
\noindent\textbf{COMPLETION: 0\%}
The calculation is an auxiliary or restricted attempt and establishes no newly verified portion of the full source target. No intrinsic criterion reconstructing a free-product-system dilation from an arbitrary noise is established.

\subsection{Sources}
\begin{itemize}
\item[\hypertarget{source:statement:1}{S1}] ULAM.AI and the UnsolvedMath Contributors, UnsolvedMath, catalogue entry. Catalogue curation under CC BY 4.0; definitions and English exposition added locally.
\url{https://huggingface.co/datasets/ulamai/UnsolvedMath}.
\item[\hypertarget{source:statement:2}{S2}] Mini-Workshop: Product Systems and Independence in Quantum Dynamics (2009), Oberwolfach Reports, DOI 10.4171/owr/2009/09, pp. 493--548; official report source and mathematical context.
\url{https://doi.org/10.4171/owr/2009/09}.
\item[\hypertarget{source:statement:3}{S3}] Official EMS Press PDF of the same Oberwolfach report; PDF pages 37--38, printed pages 529--530.
\url{https://ems.press/content/serial-article-files/46209}.
\item[\hypertarget{source:statement:4}{S4}] Skeide free product systems; reference in the saved dated catalogue literature review.
\url{https://iris.unimol.it/handle/11695/18113}.
\end{itemize}
\label{end:statement}
\end{document}
